% Ch01 WS1 -- the Day 1 scored diagnostic. 25 marks, ~30 minutes.
\documentclass{apchem-worksheet}

\apcunitword{Chapter}
\apcunit{1}
\apcunittitle{Chemical Foundations}
\apcdoctitle{Day 1 Diagnostic \textbullet\ 25 marks}
\apcsubtitle{Zumdahl \S1.3--1.10 \textbullet\ 30 minutes \textbullet\ this is scored, but it is not a test grade}

\begin{document}
\wsheader[$K = {}^\circ$C $+\ 273.15$ \textbullet\ $d = m/V$ \textbullet\
$1~\text{in} \equiv 2.54~\text{cm}$ (exact) \textbullet\
$1~\text{mi} = 5280~\text{ft}$ \textbullet\ $1~\text{ft} = 12~\text{in}$.
\textbf{Show your work} --- the point of this is to find out where you are,
not to sort you.]

\begin{apcnote}[What this is for]
This diagnostic tells your teacher which skills to reteach and which to
leave alone. It does not go in the gradebook as a test. Answer what you can,
mark anything you are guessing at, and leave blank anything you have genuinely
never seen --- a blank is far more useful information than a guess.
\end{apcnote}

\problem[]{\textbf{[6 marks]} Give the number of significant figures:}
\begin{parts}
  \item $0.00470$ \answerblank[1.6cm]{3}
  \item $6500$ \answerblank[1.6cm]{2}
  \item $6.500\times10^{3}$ \answerblank[1.6cm]{4}
  \item $2.0080$ \answerblank[1.6cm]{5}
  \item $150.$ \answerblank[1.6cm]{3}
  \item $0.0090$ \answerblank[1.6cm]{2}
\end{parts}

\problem[]{\textbf{[4 marks]} Carry out each calculation and report it to
the correct number of significant figures:}
\begin{parts}
  \item $3.24 \times 0.0150$ \workspace[1.6cm]
        \keyonly{\begin{apcanswer}$= 0.0486$; three significant figures in
        each factor, so $\mathbf{0.0486}$\end{apcanswer}}
  \item $88.4 \div 2.0$ \workspace[1.6cm]
        \keyonly{\begin{apcanswer}$= 44.2$; the limiting factor $2.0$ has
        two significant figures, so $\mathbf{44}$\end{apcanswer}}
  \item $14.31 + 2.7$ \workspace[1.6cm]
        \keyonly{\begin{apcanswer}$= 17.01$; addition, so \emph{decimal
        places} --- $2.7$ has one, giving
        $\mathbf{17.0}$\end{apcanswer}}
  \item $9.00 - 8.9$ \workspace[1.6cm]
        \keyonly{\begin{apcanswer}$= 0.10$; one decimal place, so
        $\mathbf{0.1}$. Note how much precision subtraction of close
        numbers destroys.\end{apcanswer}}
\end{parts}

\problem[]{\textbf{[4 marks]} Unit conversions. Show the conversion factors
you use.}
\begin{parts}
  \item Convert \SI{2.50}{\litre} to millilitres.
        \answerblank[3cm]{\SI{2500}{\milli\litre}}
  \item Convert \SI{350}{\milli\gram} to grams.
        \answerblank[3cm]{\SI{0.350}{\gram}}
  \item Convert \SI{0.0450}{\metre} to millimetres.
        \answerblank[3cm]{\SI{45.0}{\milli\metre}}
  \item Convert 5.00 miles to kilometres. \workspace[2.6cm]
        \keyonly{\begin{apcanswer}$5.00 \times 5280 \times 12 \times 2.54
        \div 100 \div 1000 = 8.0467\ldots$; the conversion factors are all
        exact, so three significant figures from the $5.00$:
        $\mathbf{8.05}~\text{km}$\end{apcanswer}}
\end{parts}

\problem[]{\textbf{[3 marks]} Temperature.}
\begin{parts}
  \item \SI{37.0}{\celsius} in kelvin: \answerblank[2.6cm]{310.2 K}
  \item \SI{373.15}{\kelvin} in degrees Celsius:
        \answerblank[2.6cm]{\SI{100.00}{\celsius}}
  \item Why is the kelvin scale the one used in gas-law and thermodynamic
        equations? \answerlines[2]{Because it has a true zero: zero kelvin
        is the absence of thermal energy, so quantities are directly
        proportional to $T$ in kelvin. The Celsius zero is arbitrary, and
        ratios and proportionalities fail on it.}
\end{parts}

\problem[]{\textbf{[4 marks]} Density.}
\begin{parts}
  \item A block of aluminium ($d = \SI{2.70}{\gram\per\centi\metre\cubed}$)
        has a volume of \SI{15.0}{\centi\metre\cubed}. Find its mass.
        \answerblank[3cm]{\SI{40.5}{\gram}}
  \item Ethanol has $d = \SI{0.789}{\gram\per\milli\litre}$. Find the mass
        of \SI{25.0}{\milli\litre}. \workspace[1.8cm]
        \keyonly{\begin{apcanswer}$(0.789)(25.0) =
        \mathbf{19.7}~\text{g}$\end{apcanswer}}
  \item A \SI{45.2}{\gram} metal sample raises the water level in a
        graduated cylinder from \SI{25.0}{\milli\litre} to
        \SI{41.7}{\milli\litre}. Calculate its density.
        \workspace[2.4cm]
        \keyonly{\begin{apcanswer}$V = 41.7 - 25.0 = 16.7$~mL;
        $d = 45.2/16.7 =
        \mathbf{2.71}~\si{\gram\per\milli\litre}$\end{apcanswer}}
\end{parts}

\problem[]{\textbf{[2 marks]} Precision and accuracy. A student weighs the
same object four times and records 12.47, 12.48, 12.46 and
\SI{12.48}{\gram}. The object's true mass is \SI{13.02}{\gram}.}
\begin{parts}
  \item Are the measurements precise? Accurate?
        \answerblank[6cm]{precise but not accurate}
  \item What kind of error does this indicate, and what would you do about
        it? \answerlines[2]{A systematic error --- the balance reads
        consistently low by about \SI{0.55}{\gram}. It needs
        recalibrating; taking more readings would not help.}
\end{parts}

\problem[]{\textbf{[2 marks]} Classify each as an element, a compound, a
homogeneous mixture, or a heterogeneous mixture:}
\begin{parts}
  \item Pure copper wire \answerblank[4.4cm]{element}
  \item Table salt, \ce{NaCl} \answerblank[4.4cm]{compound}
  \item Salt water \answerblank[4.4cm]{homogeneous mixture}
  \item Sand in water \answerblank[4.7cm]{heterogeneous mixture}
\end{parts}

\begin{apcnote}[Scoring guide --- for the teacher]
\textbf{Total 25 marks.} One mark per numbered part unless stated.
Award the mark only when the significant figures are also correct on
problems 2--5; that is the skill being measured.

\medskip
\textbf{22--25} --- fluent. No reteaching needed; move straight into
Unit 1. \\
\textbf{17--21} --- sound but leaky. Assign `ch01-ws2' as homework and move
on. \\
\textbf{12--16} --- the significant-figure and conversion habits are not
yet automatic. Assign `ch01-ws2', and expect to re-flag sig figs in the
first two Unit 1 worksheets. \\
\textbf{below 12} --- these skills will block progress in Unit 3
(stoichiometry) and Unit 6 (thermochemistry). Worth a dedicated float block
before Unit 1's exam.

\medskip
\textbf{Watch the pattern, not just the total.} Problems 1--2 are
significant figures, 3 is conversions, 5 is density. A student who scores
well overall but loses both marks on 2(c) and 2(d) has one specific,
fixable misconception --- applying the multiplication rule to sums --- and
does not need general remediation.
\end{apcnote}

\end{document}
